\todo{Lecture 9}

On souhaite calculer la vitesse d'un élément de fluide en fonction de sa hauteur, puis le débit volumique engendre sous ces mêmes considérations. On étudie l'écoulement laminaire dans une section avec des parois immobiles. On suppose que $\rho$ est constante, on néglige la pesanteur et $\boldsymbol{v}(\boldsymbol{r},t)=v_{x}(y)\boldsymbol{\hat{e}}_{x}$. L'équation de continuité se réduit a $\boldsymbol{\nabla v}=0$ donc
$$
\frac{ \partial u_{x} }{ \partial x } +\frac{ \partial u_{y} }{ \partial y }+\frac{ \partial u_{z} }{ \partial z } =0 
$$
et par Navier-Stokes, on trouve 
\begin{align*}
\frac{\rho D\boldsymbol{v}}{Dt} & =-\boldsymbol{\nabla}p+\eta\Delta \boldsymbol{v} \\
 & =\left( \frac{ \partial  }{ \partial t } +\boldsymbol{v\nabla} \right)\boldsymbol{v} \\
 & =\left( \frac{ \partial  }{ \partial t } +u_{x}(y)\frac{ \partial  }{ \partial x }  \right)\begin{pmatrix}
u_{x}(y) \\
0 \\
0
\end{pmatrix} =\begin{pmatrix}
0 \\
0 \\
0
\end{pmatrix} \\
\end{align*}
Donc $0=-\boldsymbol{\nabla}p+\eta\Delta \boldsymbol{v}$ et
\begin{itemize}
\item selon $y$ : $0=-\frac{ \partial p }{ \partial y }+0\Rightarrow p=f(x,z)$
\item selon $z$ : $0=-\frac{ \partial p }{ \partial z }+0\Rightarrow p=f(x)$
\item selon $x$ : $-\frac{ \partial p }{ \partial x }+\eta \frac{ \partial^{2} }{ \partial y^{2} }v_{x}(y)\Rightarrow f'(x)=\mathrm{const}\Rightarrow p(x)=-c_{1}x+c_{2}$. On a $p(0)=c_{2}=p_{\mathrm{in}}$ donc $p(\ell)=-c_{1}\ell +p_{\mathrm{in}}=p_{\mathrm{out}}$ donc $c_{1}=\frac{p_{\mathrm{in}}-p_{\mathrm{out}}}{\ell}$ et 
$$
p=p_{\mathrm{in}}-\frac{p_{\mathrm{in}}-p_{\mathrm{out}}}{\ell}x=p_{\mathrm{in}}-\frac{\Delta p}{\ell}x
$$
\end{itemize}

On veut déterminer $v_{x}(y)$. On a 
$$
\eta \frac{ \partial^{2} }{ \partial y^{2} } v_{x}=\frac{ \partial p }{ \partial x } =-c_{1}
$$
donc
$$
u_{x}=-\frac{c_{1}}{2\eta}y^{2}+ay+b
$$
qu'on peut évaluer en $\pm \frac{h}{2}$:
$$
u_{x}\left( \pm \frac{h}{2} \right)=-\frac{c_{1}}{2\eta} \frac{h^{2}}{4} \pm a \frac{h}{2}+b=0
$$
donc 
$$
a=0,\quad b=\frac{c_{1}h^{2}}{8\eta}
$$
et
$$
v_{x}(y)=\frac{c_{1}}{2\eta}\left( \frac{h^{2}}{4}-y^{2} \right)=\frac{\Delta p}{2\eta \ell}\left( \frac{h^{2}}{4}-y^{2} \right)
$$

On veut calculer le débit volumique $D$ de l'écoulement (en $\mathrm{m}^{3}\mathrm{s}^{-1}$) par une extension de l'écoulement selon $z$ de $s _{z}$ :
\begin{align*}
D & =\iint _{\mathrm{section}} \boldsymbol{v}~d\boldsymbol{S} \\
 & =\int _{-\eta /2}^{\eta /2}\int _{0}^{s _{z}} dy~dz~u_{x}(y) \\
 & = s _{z}\int_{-\eta /2}^{\eta /2} dy~ \frac{\Delta p}{2\eta \ell }\left( \frac{h^{2}}{4}-y^{2} \right) \\
 &=s _{z}\left( \frac{\Delta ph^{3}}{8\eta \ell}-\frac{\Delta p}{2\eta \ell} \frac{1}{3} y^{3}  \right)^{\eta /2}_{\eta /2} \\
 & =\frac{\Delta ps _{z} h^{3}}{12\eta \ell} 
\end{align*}
donc un debit $D$ necessite une chute de pression $\Delta p=\frac{12\eta \ell D}{s _{z} h^{3}}$
\begin{remark}
Pour un tube cylindrique $\Delta p=\frac{8\eta \ell D}{\pi R^{4}}$
\end{remark}



\chapter{Phénomènes ondulatoires}

\begin{definition}
Une onde est une perturbation d'un certain nombre de quantités physiques, qui se propagent dans l'espace, dans le cas idéale, sans déformation.
\end{definition}

\section{Onde transverse et longitudinale et équation d'onde}

\begin{definition}[Onde transverse]
Perturbation perpendiculaire a la direction de propagation de l'onde (ex. corde)
\end{definition}
\begin{definition}[Onde longitudinale]
Perturbation parallele a la direction de propagation de l'onde. (ex. ressort)
\end{definition}
\begin{remark}
Une onde transporte de l'energie/ de l'information. Pas de matiere !
\end{remark}

On peut donner une description mathematique d'une corde vibrante. Soit une corde tendue attache aux points fixes $p_{1}$ et $p_{2}$. Pour des petites vibrations, corde souple et une gravitation negligeable on peut ecrire :
$$
\frac{ \partial^{2} y_{0} }{ \partial t^{2} } =\frac{T}{\mu} \frac{ \partial^{2}y_{0} }{ \partial x^{2} }
$$
ou $\mu$ est la masse par unite de longuer de la corde, mesure en $\mathrm{kgm}^{-1}$, et $T$ est la tension de la corde, qu'on mesure en $\mathrm{N}=\mathrm{kgms}^{-2}$. Prenons $y_{0}(x,t)=f(x-ct)$ comme une solution. Alors
\begin{align*}
\frac{ \partial  }{ \partial t } f(x-ct) & =f'(x-ct) \frac{ \partial  }{ \partial t } (x-ct) \\
 & =cf'(x-ct) \\
\Rightarrow \frac{ \partial^{2} }{ \partial t^{2} } f(x-ct) & =c^{2}f''(x-ct)
\end{align*}
et de meme
$$
\frac{ \partial^{2} }{ \partial x^{2} } f(x-ct)=f''(x-ct)
$$
donc
$$
\frac{ \partial^{2}y_{0} }{ \partial t^{2} } =c^{2}f''(x-ct)\overset{?}=\frac{T}{\mu}\frac{ \partial^{2}y_{0} }{ \partial x^{2} } =\frac{T}{\mu}f''(x-ct)
$$
alors on a que
$$
c=\sqrt{ \frac{T}{\mu} }
$$

Soient $f$ et $g$ deux solutions d'une equation differentielle lineaire en $y_{0}$ alors $f+g$ est aussi une solution. Donc la solution generale est de la forme
$$
y_{0}(x,t)=f(x-ct)+g(x+ct)
$$


\section{Ondes sinusoidales}

On a un cas particulier des ondes de la forme
$$
y_{0}(x,t)=A\cos (\omega t-kx+\varphi)=A\cos\left( -k\left( x-\frac{\omega}{k}t \right)+\varphi \right)
$$
qui est une fonction du type $f(x-ct)$ avec $c=\frac{\omega}{k}$. Dans cette fonction on a $A$ l'amplitude, $k$ le nombre d'onde ($\mathrm{rad~m}^{-1}$), $\omega$ la pulsation ou frequence angulaire ($\mathrm{rad~s}^{-1}$) et $\varphi$ le deplacement ($\mathrm{rad}$).
\begin{definition}[Longueur d'onde]
Periode spatiale de l'onde c-a-d distance entre oscillations. On la note
$$
\lambda=\frac{2\pi}{k},\quad[\lambda]=\mathrm{m}
$$
\end{definition}
\begin{definition}[Periode]
Periode temporelle de l'onde c-a-d temps entre oscillations. On la note
$$
T=\frac{2\pi}{\omega},\quad [\omega]=\mathrm{s}
$$
\end{definition}
\begin{definition}[Frequence]
Le nombre d'oscillations par unite de temps. On la note
$$
\nu=\frac{1}{T}=\frac{\omega}{2\pi},\quad [\nu]=\mathrm{Hz}=\mathrm{s}^{-1}
$$
\end{definition}
\begin{remark}
Souvent, on utilise la notation complexe
$$
\tilde{y}_{0}(x,t)=\tilde{A}e^{i(\omega t-kx)},\quad \tilde{A}=Ae^{i\varphi}
$$
en rappelent que la partie physique est la partie reele
$$
\mathrm{Re}(Ae^{i(\omega t-kx+\varphi)})=A\cos(\omega t-kx+\varphi)
$$
\end{remark}